\subsection{What would close it, computed rather than listed}
\label{ssec:closure}

A list of open problems invites a reading it does not deserve, that each entry
is a step and that the steps together are a path. The implications proved
above, together with those of \Cref{ssec:bullseye}, decide the question, so
we settle it by computation instead of by description. The logical skeleton
of the paper and of the literature it quotes is finite: a set of statements,
each of them proved here, quoted from the literature, or open, and a set of
inference rules, each of them one theorem. Consequences can then be taken.

Such a computation is only as strong as its rule set. A set of statements
the rules do not carry to the conjecture is \emph{insufficient} in the sense
of the computation, and that means no more than that no rule derives the
conjecture from it. An implication that is true and left out makes the
computation report a set as necessary when it is not. One such implication
is elementary, and it decides the shape of the answer, so we prove it first.

\begin{proposition}[The conjecture for the varieties that are not abelian
is the conjecture]\label{prop:f3ishc}
Let $A$ be a complex abelian variety of dimension $g\ge1$ and let
$X=A\times\PP^{1}$. Then $X$ is a smooth complex projective variety that is
not an abelian variety, and if every Hodge class of degree $2p$ on $X$ is
algebraic, then so is every Hodge class of degree $2p$ on $A$. Consequently
the statement
\begin{quote}
the Hodge conjecture holds for every smooth complex projective variety that
is not an abelian variety
\end{quote}
implies the Hodge conjecture for every abelian variety, and it is equivalent
to $\mathbf{HC}$.
\end{proposition}

\begin{proof}
$X$ is smooth and projective of dimension $g+1$. By the K\"unneth formula
$h^{1,0}(X)=h^{1,0}(A)+h^{1,0}(\PP^{1})=g$, while an abelian variety of
dimension $g+1$ has $h^{1,0}=g+1$; so $X$ is not an abelian variety. (Equally,
$X$ contains the rational curves $\{a\}\times\PP^{1}$, and every morphism
from $\PP^{1}$ to a complex torus is constant, since it lifts to the
universal cover $\CC^{g+1}$ \cite[Chapter~1]{BL04}.)

Let $\mathrm{pr}\colon X\to A$ be the projection and $D=A\times\{0\}$, a
smooth divisor that $\mathrm{pr}$ maps isomorphically onto $A$, so that
$\mathrm{pr}_{*}[D]=1$ in $H^{0}(A,\QQ)$. Let $\alpha$ be a Hodge class of
degree $2p$ on $A$. Pull-back along a morphism is a morphism of Hodge
structures, so $\mathrm{pr}^{*}\alpha$ is a Hodge class of degree $2p$ on
$X$, and by hypothesis $\mathrm{pr}^{*}\alpha=[Z]$ for an algebraic cycle
$Z$ on $X$ with rational coefficients. The projection formula gives
\[
  \mathrm{pr}_{*}\bigl([Z]\cup[D]\bigr)
  =\mathrm{pr}_{*}\bigl(\mathrm{pr}^{*}\alpha\cup[D]\bigr)
  =\alpha\cup\mathrm{pr}_{*}[D]=\alpha .
\]
The cycle class map is compatible with intersection products and with proper
push-forward \cite[Chapter~19]{Ful98}, so the left side
is the class of the cycle $\mathrm{pr}_{*}(Z\cdot D)$ and $\alpha$ is
algebraic. A point, the abelian variety of dimension zero, satisfies the
conjecture trivially.

So the displayed statement gives the conjecture for every abelian variety,
and with it the conjecture for every smooth complex projective variety. The
converse is immediate: the displayed statement is a special case of
$\mathbf{HC}$.
\end{proof}

Since the conjecture for the varieties that are not abelian is the whole
conjecture, the statement that belongs on a frontier is weaker: it asks for
algebraicity only modulo what abelian varieties supply.

\begin{definition}[The conjecture modulo abelian varieties]
\label{def:f3prime}
For a smooth complex projective variety $X$ and $p\ge0$ let
$\mathrm{Ab}^{p}(X)\subseteq H^{2p}(X,\QQ)$ be the span of the classes
$\Gamma_{*}\gamma=\mathrm{pr}_{X*}\bigl([\Gamma]\cup\mathrm{pr}_{B}^{*}\gamma\bigr)$
of degree $2p$, where $B$ runs over the complex abelian varieties, the point
included, $\gamma$ over the Hodge classes on $B$, and $\Gamma$ over the
algebraic cycles on $B\times X$ with rational coefficients. The statement
\textup{(F3$'$)} is: for every $X$ and every $p$,
\[
  \Hdg^{p}(X)=\mathrm{Ab}^{p}(X).
\]
\end{definition}

Each $\Gamma_{*}\gamma$ is a Hodge class, since cup product with an
algebraic class and push-forward preserve Hodge classes, so
$\mathrm{Ab}^{p}(X)\subseteq\Hdg^{p}(X)$; and taking $B$ a point shows that
$\mathrm{Ab}^{p}(X)$ contains every algebraic class. So \textup{(F3$'$)} says
that every Hodge class is algebraic modulo the images of Hodge classes of
abelian varieties under algebraic correspondences, which is the form of (F3)
that \Cref{rem:f3collapse} asks for.

\begin{proposition}[What is known of \textup{(F3$'$)}]\label{prop:f3prime}
\leavevmode
\begin{enumerate}[label=\textup{(\roman*)},leftmargin=2.6em]
\item $\mathbf{HC}$ implies \textup{(F3$'$)}, and \textup{(F3$'$)} together
  with the conjecture for abelian varieties implies $\mathbf{HC}$.
\item \textup{(F3$'$)} holds on every abelian variety and on every product
  $A\times\PP^{1}$ with $A$ abelian; the argument of \Cref{prop:f3ishc} gives
  nothing for it.
\item Let $\mathcal{A}$ be the class of smooth complex projective varieties
  $X$ for which $H^{*}(X,\QQ)$ is spanned by classes $\Gamma_{*}y$ with $B$
  abelian, $\Gamma$ an algebraic cycle on $B\times X$ and
  $y\in H^{*}(B,\QQ)$. Every $X\in\mathcal{A}$ satisfies \textup{(F3$'$)}.
  The class $\mathcal{A}$ contains every abelian variety and every smooth
  projective curve, and it is closed under products and under images of
  surjective morphisms. So it contains every product of curves, every
  symmetric product of a curve, and every smooth projective variety onto
  which a product of curves or an abelian variety maps surjectively.
\item Every $X$ of dimension $n$ satisfies \textup{(F3$'$)} in the degrees
  $2p$ with $p\in\{0,1,n-1,n\}$, so every $X$ of dimension at most three
  satisfies it in every degree.
\end{enumerate}
\end{proposition}

\begin{proof}
(i) If $\mathbf{HC}$ holds, then $\Hdg^{p}(X)$ consists of algebraic classes,
which lie in $\mathrm{Ab}^{p}(X)$. Conversely, if the conjecture holds for
abelian varieties, every $\gamma$ in \Cref{def:f3prime} is algebraic, so
every $\Gamma_{*}\gamma$ is, since composition with a correspondence carries
algebraic classes to algebraic classes \cite[Chapter~16]{Ful98}; then
\textup{(F3$'$)} says that every Hodge class is algebraic.

(ii) On an abelian variety $A$ take $B=A$ and $\Gamma$ the diagonal. On
$X=A\times\PP^{1}$ the K\"unneth formula gives
$H^{*}(X,\QQ)=\mathrm{pr}^{*}H^{*}(A,\QQ)\oplus
\mathrm{pr}^{*}H^{*}(A,\QQ)\cup\eta$, with $\eta$ the class of a point of
$\PP^{1}$ pulled back to $X$, which is algebraic. The first summand is
$\Gamma_{*}H^{*}(A,\QQ)$ for $\Gamma$ the transpose of the graph of
$\mathrm{pr}$, and the second is $\Gamma'_{*}H^{*}(A,\QQ)$ for
$\Gamma'=\Gamma\cdot\mathrm{pr}_{X}^{*}\eta$. So $X\in\mathcal{A}$, and
(iii) applies.

(iii) Let $X\in\mathcal{A}$ and fix $p$. Since $H^{2p}(X,\QQ)$ is finite
dimensional, finitely many classes $\Gamma_{j*}y_{j}$ span it, and we may
take each $\Gamma_{j}$ of pure codimension $c_{j}$ on $B_{j}\times X$ and each
$y_{j}$ of pure degree $k_{j}$, with $k_{j}+2c_{j}-2\dim B_{j}=2p$. Then
$\Gamma_{j*}\colon H^{k_{j}}(B_{j},\QQ)\to H^{2p}(X,\QQ)$ is a morphism of
Hodge structures up to a Tate twist, and the sum of these maps is a
surjective morphism of polarisable Hodge structures onto $H^{2p}(X,\QQ)$.
Polarisable Hodge structures are semisimple, so its kernel has a complement
that is a sub-Hodge structure and maps isomorphically onto $H^{2p}(X,\QQ)$;
hence every Hodge class of degree $2p$ is the image of a Hodge class, that is
$\sum_{j}\Gamma_{j*}\gamma_{j}$ with each $\gamma_{j}$ a Hodge class on
$B_{j}$. So $\Hdg^{p}(X)=\mathrm{Ab}^{p}(X)$.

An abelian variety lies in $\mathcal{A}$ through its diagonal. A curve $C$
of genus $g\ge1$ lies in $\mathcal{A}$ through an Abel--Jacobi embedding
$\iota\colon C\to J(C)$: the transpose of its graph induces $\iota^{*}$,
which is an isomorphism on $H^{0}$ and $H^{1}$ and sends the class of the
theta divisor to $g$ times the class of a point, so it is surjective; the
projective line is the image of an elliptic curve under a double cover, and
is covered by the last closure property. For products, if $H^{*}(X_{i},\QQ)$
is spanned by the $\Gamma_{i*}y$, then by the K\"unneth formula
$H^{*}(X_{1}\times X_{2},\QQ)$ is spanned by the
$(\Gamma_{1}\times\Gamma_{2})_{*}(y_{1}\otimes y_{2})$, with
$\Gamma_{1}\times\Gamma_{2}$ an algebraic cycle on
$(B_{1}\times B_{2})\times(X_{1}\times X_{2})$. For images, let
$f\colon Y\to X$ be a surjective morphism with $Y\in\mathcal{A}$, let
$r=\dim Y-\dim X$ and let $h$ be an ample class on $Y$. The correspondence
$\Psi=[\Gamma_{f}]\cdot\mathrm{pr}_{Y}^{*}h^{r}$ on $Y\times X$ induces
$\Psi_{*}y=f_{*}(y\cup h^{r})$, and
$\Psi_{*}f^{*}x=x\cup f_{*}h^{r}=c\,x$ with $c>0$ the degree of $h^{r}$ on a
general fibre. So $\Psi_{*}$ is surjective, and $H^{*}(X,\QQ)$ is spanned by
the classes $(\Psi\circ\Gamma)_{*}y$, with $\Psi\circ\Gamma$ algebraic
\cite[Chapter~16]{Ful98}. A symmetric product of a curve is the image of a
power of it.

(iv) Every class of degree $0$ or $2n$ is algebraic, and every Hodge class of
degree $2$ is algebraic by the Lefschetz theorem on $(1,1)$ classes. For
$n\ge2$, the hard Lefschetz theorem makes cup product with $h^{n-2}$, for $h$
an ample class, an isomorphism $H^{2}(X,\QQ)\to H^{2n-2}(X,\QQ)$ of Hodge
structures up to a Tate twist, so every Hodge class of degree $2n-2$ is
$h^{n-2}$ times a Hodge class of degree $2$, and is algebraic. Algebraic
classes lie in $\mathrm{Ab}^{p}(X)$. When $n\le3$ every $p$ between $0$ and
$n$ is one of $0,1,n-1,n$.
\end{proof}

\begin{remark}[Why \textup{(F3$'$)} is not closed here]\label{rem:f3primeopen}
\textup{(F3$'$)} is the statement that the Hodge conjecture reduces to
abelian varieties along algebraic correspondences, and
\Cref{prop:f3prime} records what is known of such a reduction: the
varieties of $\mathcal{A}$, whose whole cohomology is reached from abelian
varieties by algebraic correspondences, and the degrees in which the
conjecture itself is known. The first degree in which it is not known in
general is degree four on fourfolds, for fourfolds not known to lie in
$\mathcal{A}$. The square $S\times S$ of a K3 surface is the standard
example. Its Hodge classes in $H^{2}(S)\otimes H^{2}(S)$ include the classes
that induce the Hodge endomorphisms of the transcendental lattice $T(S)$, and
$S\times S$ lies in $\mathcal{A}$ as soon as the Kuga--Satake embedding of
$T(S)$ is induced by an algebraic cycle, the Kuga--Satake Hodge conjecture for
$S$, which is not known in general: the transpose of such a cycle maps the
cohomology of the Kuga--Satake variety onto $T(S)$, and the rest of
$H^{*}(S,\QQ)$ is algebraic. The rational Hodge isometries are algebraic
\cite{Bus19}, but a real multiplication is not an isometry, and
\Cref{rem:mumfordks} meets exactly this obstacle for the surfaces
$S_{\lambda}$.

For varieties whose cohomology is not of abelian type the position is
starker. By \Cref{prop:noabeliantype} no correspondence moves the $H^{3}$ of a
very general quintic threefold into the cohomology of an abelian variety, so
such a variety does not lie in $\mathcal{A}$. That is no obstruction to
\textup{(F3$'$)} itself, since a Hodge class spans a Hodge structure of Tate
type, which is of abelian type, and nothing in Hodge theory forbids it from
lying in $\mathrm{Ab}^{p}(X)$. Nor is it a route: to put a Hodge class into
$\mathrm{Ab}^{p}(X)$ one must construct an algebraic cycle on $B\times X$
whose class carries it, which is the kind of construction the Hodge
conjecture asks for, now on $B\times X$ instead of on $X$. The methods of
this paper construct cycles on abelian varieties of Weil type and on nothing
else, so they give no case of \textup{(F3$'$)} outside $\mathcal{A}$. We
therefore leave it open, and \Cref{thm:closure} treats it as an open
statement.
\end{remark}

\begin{setupenv}[The rule set]\label{setup:rules}
Item (XXXIII) of \Cref{sec:verification} carries forty-two statements and
twenty-two rules; the propagation statement (P2) is carried family by
family, as (P2) for the split and for the other families of the imaginary
quadratic fields and of the fields of degree at least four. Each statement is marked proved in this paper, quoted from the
literature, or open. Each rule is a finite set of premises together with one
conclusion, and carries the label of the theorem of this paper that proves it;
\Cref{tab:rules} lists the rules whose conclusion is derived rather than
assumed. For a set $T$ of statements write $\Cn(T)$ for the least set that
contains $T$ and is closed under the rules, which exists because each rule has
a single conclusion and the consequence operator is therefore monotone. Write
$\mathbf{T}$ for the set of statements marked proved or quoted, and
$\mathbf{HC}$ for the Hodge conjecture for every smooth complex projective
variety. Call a set $S$ of open statements \emph{sufficient} if
$\mathbf{HC}\in\Cn(\mathbf{T}\cup S)$, and \emph{minimal} if no proper subset
of $S$ is sufficient.
\end{setupenv}

\begin{table}[htbp]
\centering
\small
\renewcommand{\arraystretch}{1.14}
\begin{tabular}{@{}>{\raggedright\arraybackslash}p{0.40\linewidth}>{\raggedright\arraybackslash}p{0.27\linewidth}l@{}}
\toprule
premises & conclusion & by \\
\midrule
the conjecture for abelian varieties, (F3)
  & $\mathbf{HC}$ & \Cref{ssec:closuremap} \\
(F3)
  & abelian varieties & \Cref{prop:f3ishc} \\
the conjecture for abelian varieties, (F3$'$)
  & $\mathbf{HC}$ & \Cref{prop:f3prime} \\
the Weil classes of every CM field, (F2)
  & abelian varieties & \Cref{ssec:closuremap} \\
$K$ imaginary quadratic, the fields of degree at least four
  & every CM field & \Cref{app:albert} \\
$\delta=\delta_{0}$, $\delta\ne\delta_{0}$
  & every discriminant & \Cref{prop:separate} \\
a base point, the reduction, the orbit, (P2) for the family
  & that family & \Cref{thm:semiregclosure} \\
the same three for a CM field of degree at least four, (P2) for the family
  & that family & \Cref{thm:cmpropagation} \\
its split families, its other families
  & the fields of degree at least four & \Cref{prop:cmweil} \\
$\delta=\delta_{0}$
  & $\delta\ne\delta_{0}$ & \Cref{prop:descent} \\
the split families of the fields of degree at least four
  & their other families; $\delta=\delta_{0}$ & \Cref{prop:descent} \\
a secant object in every dimension, Question 11.4
  & $\delta=\delta_{0}$ & \Cref{prop:splitgeom} \\
a Cohen--Macaulay support, smooth or singular
  & one secant object & \Cref{cor:disjointroutes} \\
one secant object, Question 11.4
  & $\mathbf{W}(K,4,\delta_{0})$ & \Cref{cor:markmancondition} \\
\midrule
(L), (M)
  & $\mathbf{HC}$ & \Cref{prop:lefmot} \\
(L), the deformation theorem for motivated classes
  & (V) & \Cref{prop:lefvhc} \\
(V), the accessibility of Hodge classes on abelian varieties
  & abelian varieties & \Cref{thm:vhcab} \\
\bottomrule
\end{tabular}
\caption{The rules of \Cref{setup:rules} whose conclusion is derived, with the
theorem that proves each; the rows for (P2) and for the Cohen--Macaulay
support each stand for two rules, the row for descent over the fields of
degree at least four for two, descent within a field and scalar extension
down to the imaginary quadratic fields, the second row is \Cref{prop:f3ishc},
which an earlier version of the set omitted, the third is
\Cref{prop:f3prime}(i) for the weaker statement (F3$'$) of
\Cref{def:f3prime}, the last three rows are the rules of
\Cref{ssec:bullseye}, one of whose conclusions, (V), is itself open, and the
one remaining rule records what the literature supplies on Weil classes. The
bounded criterion (P1) is absent: by
\Cref{thm:p1equivalent} it is equivalent to its own conclusion, so it is a
restatement and not a premise. Item (XXXIII) of
\Cref{sec:verification} carries the whole set as data and checks that every
label here is a label of this paper.}
\label{tab:rules}
\end{table}

\begin{theorem}[The closure of what is proved here]\label{thm:closure}
With the notation of \Cref{setup:rules}:
\begin{enumerate}[label=\textup{(\roman*)},leftmargin=2.6em]
\item $\mathbf{HC}\notin\Cn(\mathbf{T})$. Nothing in this paper, combined with the
  theorems it quotes, implies the Hodge conjecture.
\item There are exactly five minimal sufficient sets:
  \[
    \{\mathrm{F3}\},\quad\{\mathrm{L},\mathrm{M}\},\quad
    \{\mathrm{L},\mathrm{F3}'\},\quad\{\mathrm{V},\mathrm{F3}'\},\quad
    \{\mathrm{P2}_{\mathrm{split}},\mathrm{F2},\mathrm{F3}'\},
  \]
  with \textup{(L)}, \textup{(M)}, \textup{(V)} as in \Cref{def:lmv},
  \textup{(F2)}, \textup{(F3)} as in \Cref{cor:fourremain}, (F3$'$) as in
  \Cref{def:f3prime}, and $\mathrm{P2}_{\mathrm{split}}$ the propagation
  statement \textup{(P2)}, read in the corrected form of \Cref{rem:p2prime},
  for the split families of the CM fields of degree at least four. No
  minimal set contains \textup{(P2)} for an imaginary quadratic field or for
  a family of nontrivial discriminant: \Cref{prop:descent} derives those.
  The single open statement (F3) is sufficient, and it is equivalent to
  $\mathbf{HC}$ (\Cref{prop:f3ishc}); no other single open statement is.
  The sets $\{\mathrm{L},\mathrm{F3}\}$, $\{\mathrm{V},\mathrm{F3}\}$ and
  $\{\mathrm{P2}_{\mathrm{split}},\mathrm{F2},\mathrm{F3}\}$ are sufficient and
  not minimal. The last of the five is the only one with three elements and the
  only one whose derivation of $\mathbf{HC}$ uses a statement proved in this
  paper.
\item In each of the five, every element is necessary: removing any one
  leaves $\mathbf{HC}$ outside the closure.
\item The secant route lies in no minimal sufficient set. Granting both of its
  open demands, a secant object satisfying \eqref{eq:heisenberg} in every
  dimension and an affirmative answer to \cite[Question 11.4]{Mar25b}, puts
  $\mathbf{W}(K,n,\delta_{0})$ in the closure, and with it, by
  \Cref{prop:descent}, $\mathbf{W}(K,n,\delta)$ for every $\delta$; it leaves
  the Weil classes of the CM fields of degree at least four outside.
\item No statement proved in this paper rests, at any depth, on an open one,
  and the rule set is acyclic.
\item Every element of a minimal sufficient set other than
  $\mathrm{P2}_{\mathrm{split}}$ is a consequence of $\mathbf{HC}$, and $\{\mathrm{F3}\}$ and
  $\{\mathrm{L},\mathrm{M}\}$ are each equivalent to it; no implication from
  $\mathbf{HC}$ to \textup{(P2)} is known.
\item An earlier version of this theorem, over a rule set without
  \Cref{prop:descent}, stated in (iv) that the secant route stops at the
  trivial discriminant, and listed \textup{(P2)} for every family in the
  fifth set. Both came from the omission of the rules of
  \Cref{prop:descent}, and the five minimal sets are otherwise unchanged.
\end{enumerate}
\end{theorem}

\begin{proof}
The consequence operator is monotone, so $\Cn$ is computed by iterating the
one step map to a fixed point, and the set of subsets of the sixteen open
statements is finite. Item (XXXIII) of
\Cref{sec:verification} carries out the iteration and the exhaustive search
over all $2^{16}$ subsets of the sixteen open statements in exact arithmetic,
and Section 24 of the Lean certificate repeats (i), (ii), (iii) and (iv) in
the kernel, over the same rule set under the numbering recorded there,
checking in particular all sixteen single open statements and all one
hundred and twenty pairs. The content of the statement is the rule set itself, and every rule is
one of the theorems named in \Cref{tab:rules}. The sufficiency of (F3)
alone is the chain $\mathrm{(F3)}\Rightarrow$ abelian varieties by
\Cref{prop:f3ishc}, then abelian varieties $\wedge\,\mathrm{(F3)}
\Rightarrow\mathbf{HC}$; every sufficient set that contains (F3) and has
another element is therefore not minimal, and a sufficient set without
(F3) must reach $\mathbf{HC}$ through one of the two other rules with that
conclusion, \Cref{prop:lefmot} or \Cref{prop:f3prime}(i), the second of
which needs the conjecture for abelian varieties, reached by \textup{(V)},
by \textup{(L)} through \textup{(V)}, or by (P2) with (F2). For (vi):
\textup{(L)} follows from $\mathbf{HC}$ by \Cref{thm:equivalence},
\textup{(M)} by \Cref{prop:lefmot}, \textup{(V)} by \Cref{prop:lefvhc},
(F3$'$) by \Cref{prop:f3prime}(i), and (F2) and (F3) are special cases of
$\mathbf{HC}$.

In (iv), every base point the secant route produces is a member
$X\times\widehat{X}$, whose discriminant is trivial whatever $\beta$ is chosen
(\Cref{prop:splitgeom}(ii)); the deformation argument stays inside the family
of trivial discriminant, and it is \Cref{prop:descent}(i), not the
deformation, that carries the result to the other discriminants. It has no
base point over a field of degree at least four, so it reaches none of them.
\end{proof}

\begin{corollary}[The three statements, one of which is the
conjecture]\label{cor:fourremain}
The Hodge conjecture follows from the results of this paper together with the
following three statements. No rule of \Cref{setup:rules} derives it from
(F1) and (F2) together, and it follows from (F3) alone, which is equivalent
to it
(\Cref{prop:f3ishc}); so the list is not a reduction of the conjecture to
three smaller problems. With (F3) replaced by the weaker statement
\textup{(F3$'$)} of \Cref{def:f3prime}, the conjecture modulo abelian
varieties, the three statements (F1), (F2), \textup{(F3$'$)} still suffice,
and no rule of \Cref{setup:rules} derives the conjecture from a proper
subset of them.
\begin{enumerate}[label=\textup{(F\arabic*)},leftmargin=3.0em,itemsep=4pt]
\item \emph{One propagation statement, for every CM field.} \textup{(P2)},
  in the corrected form of \Cref{rem:p2prime}: at one base point of each
  Weil family of each CM field some perfect complex whose Chern character is
  a nonzero multiple of a rational Weil class plus a polynomial in the
  polarisation has injective semiregularity map. Its first form, with Chern
  character exactly a multiple of the Weil class, is false for every
  imaginary quadratic field and every $n\ge2$ (\Cref{thm:p2false}): such a
  complex would have to deform to non-algebraic Weil tori, whose coherent
  sheaves have no Chern character in degrees $1$ to $2n-1$. The corrected form
  closes
  $\mathbf{W}(K,n,\delta)$
  for every imaginary quadratic $K$, every $n$ and every $\delta$ at once, by
  \Cref{thm:semiregclosure}, and
  $\mathbf{W}(F,n,\delta)$ for every CM field $F$ of higher degree, by
  \Cref{thm:cmpropagation}, because those theorems put $\Sigma=\cD$, so the
  Weil class is algebraic at every member of the family and not only at the
  general one, and because \Cref{thm:everyfamily} and \Cref{thm:cmbasepoint}
  supply the base point in every one of those families. By
  \Cref{prop:descent} it is enough to have it for the split families of the
  CM fields of degree at least four, in infinitely many dimensions.
  \Cref{thm:p2forces}, \Cref{thm:p2support}, \Cref{thm:nosum},
  \Cref{thm:p2numerical} and \Cref{thm:p2endo} describe the object of the
  first form and are vacuous by \Cref{thm:p2false}; for the corrected form
  \Cref{thm:p2primenumber} computes the number that $\dim\Ext^{2}(E,E)$ has
  to reach, $7n^{2}-2n$ for a general shape. For one family its
  conclusion is equivalent to the Lefschetz standard conjecture for one
  variety (\Cref{cor:weilequiv}). The bounded criterion \textup{(P1)} is not a
  second form of the same demand: by
  \Cref{thm:p1equivalent} it is equivalent to the conclusion it was to imply,
  and \Cref{thm:p1forces} shows that the one cycle it could have been tested
  on, the transport of the base cycle, does not satisfy it.
\item \emph{The classes beyond the Weil lines.} Every Hodge class on an
  abelian variety that lies outside the subring generated by divisor classes
  and by Weil classes of quotients is algebraic. Through dimension five there
  is no such class \cite{MZ99}; in dimension eight there are
  (\Cref{prop:mumfordproduct}), so this is not a reduction of the conjecture
  for abelian varieties to Weil classes, which is false as a statement about
  Hodge rings: it is the conjecture for those classes. Its smallest case is
  brought down to one Kuga--Satake class in \Cref{thm:mumfordks}; it holds at
  every CM point (\Cref{thm:mumfordcm}), and on all members of a Mumford
  family it is equivalent to the Lefschetz standard conjecture for one
  ninefold (\Cref{cor:mumfordequiv}). Its classes are rigid along the base
  curve (\Cref{thm:mumfordrigid}), and one complex at one CM point with
  $\Ext^{2}$ of dimension $119$ would close it (\Cref{thm:mumfordnumerical}),
  as would a bound on the cycles that represent them at the closed points of
  the reductions of the curve modulo primes (\Cref{thm:modpdensity}).
\item \emph{The varieties that are not abelian.} The Hodge conjecture holds
  for every smooth complex projective variety that is not an abelian variety.
  This is the Hodge conjecture itself. The varieties it covers include
  $A\times\PP^{1}$ for every abelian variety $A$, and the conjecture for
  $A\times\PP^{1}$ gives it for $A$ (\Cref{prop:f3ishc}); so (F3) implies
  the conjecture for abelian varieties, hence (F1) and (F2), and hence
  $\mathbf{HC}$. Nor is it a reduction to abelian varieties in the other
  direction: no realisation of the cohomology of an arbitrary variety inside
  that of abelian varieties exists (\Cref{prop:noabeliantype}), and
  \Cref{thm:hyperplane} and \Cref{thm:returntrip} close the two routes
  through hyperplane sections that would be the natural candidates for one.
\end{enumerate}
\end{corollary}

\begin{proof}
The three together are sufficient by the derivation of \Cref{thm:closure},
and (F3) alone by \Cref{prop:f3ishc}. That (F1) and (F2) together are not
sufficient in the rule set is \Cref{thm:closure}(ii), since every minimal
sufficient set contains (F3), \textup{(F3$'$)} or \textup{(M)}; the
statement about (F1), (F2), \textup{(F3$'$)} is \Cref{thm:closure}(ii) and
(iii) for the last of the five minimal sets. The Weil classes of the CM
fields of degree at least four are not a fourth statement: they are derived
from (F1) by \Cref{thm:cmpropagation}. Nor is the bounded criterion a
premise, since by \Cref{thm:p1equivalent} it is a restatement of its own
conclusion.
\end{proof}

\begin{remark}[How the earlier count went wrong, and the wording of (F2)]
\label{rem:f3collapse}
An earlier version of this corollary said that the conjecture follows from
no proper subset of the three statements. That was a statement about a rule
set that used (F3) only through $\text{(abelian varieties)}\wedge
\text{(F3)}\Rightarrow\mathbf{HC}$ and did not contain the elementary rule
$\text{(F3)}\Rightarrow\text{(abelian varieties)}$ of \Cref{prop:f3ishc}.
The rule set can show that a set suffices; it cannot show that a set is
needed, except relative to the implications written into it.

The wording of (F2) has the same feature, one degree at a time. Let
$R(A)\subseteq H^{*}(A,\QQ)$ be the subring generated by divisor classes and
Weil classes of quotients. Suppose some Hodge class $\beta$ of degree $2p$ on
$A$ lies outside $R(A)$, and let $\alpha$ be any Hodge class of degree $2p$.
If $\alpha\notin R(A)$, (F2) makes it algebraic. If $\alpha\in R(A)$, then
$\alpha+\beta\notin R(A)$, because $R(A)$ is closed under subtraction, so
(F2) makes $\alpha+\beta$ and $\beta$ algebraic, and with them $\alpha$.
So in every degree in which (F2) says anything it is the whole conjecture for
that degree: on the square of a Mumford fourfold it gives every Hodge class of
$H^{4}$, the six products of divisor classes included. Whether (F2) alone
gives the conjecture for every abelian variety is not decided here, and
nothing depends on it: the rule set does not use such an implication, and
adding it would not change \Cref{thm:closure}(ii), since (F3) alone is
already sufficient.

The form of (F3) that is not the conjecture itself asks only for
algebraicity modulo classes coming from abelian varieties. It is
\textup{(F3$'$)} of \Cref{def:f3prime}; \Cref{prop:f3prime} says what is
known of it, and \Cref{rem:f3primeopen} why it is not closed.
\end{remark}

\begin{figure}[tbp]
\centering
  \includegraphics[width=0.92\textwidth]{fig_frontier}
\caption{The rule set of \Cref{setup:rules} with the rules of
\Cref{ssec:bullseye}, drawn as a hypergraph. A rule with several premises is
a junction marked $\wedge$, with arrows from its premises and one arrow to its
conclusion; each rule carries the result that proves it, and the junction
marked \S16.5 is the closure map of \Cref{ssec:closuremap}. Clay marks an
open statement, grass what this paper proves, indigo a derived statement,
and the tag $\Leftarrow\mathbf{HC}$ an open statement that the conjecture
implies: every clay node except (P2). The arrow from (F3) to the conjecture
is \Cref{prop:f3ishc}, and the junction above the conjecture for abelian
varieties is \Cref{prop:f3prime}(i). The five minimal sufficient sets of
\Cref{thm:closure}(ii) are drawn as nested ribbons under the arrows they
use, widest first, and listed in the cards below: (F3) alone; (L) with (M);
(L) through (V) with \textup{(F3$'$)}; (V) with \textup{(F3$'$)}; and, on
the left, (P2) with (F2) and \textup{(F3$'$)}. A shared arrow shows the
bands of every route through it. The secant route, which reaches the
imaginary quadratic families only, is drawn in \Cref{fig:closure}.}
\label{fig:frontier}
\end{figure}

\begin{remark}[The target, and what it would take]\label{rem:bullseye}
With the rules of \Cref{ssec:bullseye}, \Cref{prop:f3ishc} and
\Cref{prop:f3prime} in place, \Cref{thm:closure} changes the reading of
\Cref{cor:fourremain} in three ways. First, the route through this paper is
a route only with (F3) replaced by \textup{(F3$'$)}. As stated, its last
statement (F3) is the conjecture (\Cref{prop:f3ishc}), so (F1) and (F2), and
everything this paper proves about Weil classes, drop out of every minimal
derivation that uses (F3). With \textup{(F3$'$)} in its place the route
$\{\mathrm{P2}_{\mathrm{split}},\mathrm{F2},\mathrm{F3}'\}$ is minimal, and it is not the
shortest: (F1) and (F2), which between them give the conjecture for abelian
varieties along this route, are replaced on the routes of the literature by
the single statement \textup{(V)}, or by \textup{(L)}, which implies it
(\Cref{thm:vhcab}). Second, every element of every minimal sufficient set
except $\mathrm{P2}_{\mathrm{split}}$ is a consequence of the conjecture
(\Cref{thm:closure}(vi)), and
two of the sets, $\{\mathrm{F3}\}$ and $\{\mathrm{L},\mathrm{M}\}$, are
each equivalent to it. (P2) is a statement about semiregularity maps, and a
proof of the conjecture need not prove it; in its first form, with a Chern
character that is exactly a Weil class, it is false (\Cref{thm:p2false}), and
the route runs through the corrected form of \Cref{rem:p2prime}. Third, every minimal sufficient set
contains (F3), \textup{(F3$'$)} or \textup{(M)}: the varieties that are not
abelian are reached only through the conjecture for them, which is the whole
conjecture, through the conjecture modulo abelian varieties, which is open
(\Cref{rem:f3primeopen}), or through motivated classes.

So the target at the centre is \textup{(L)}. It lies in two of the five
minimal sets, it gives \textup{(V)} and with it the whole abelian half, the
conjecture implies it, and with \textup{(M)} it is the conjecture. For each
$X$ it asks for one explicit Hodge class on $X\times X$ to be algebraic, the
class inducing $\Lambda$. It is open, as are \textup{(M)}, \textup{(V)}
and \textup{(F3$'$)}, and this paper proves none of them in general:
\Cref{cor:mumfordlefschetz} adds the total spaces of Mumford families to the
known cases of \textup{(L)}, \Cref{prop:f3prime} records the known cases of
\textup{(F3$'$)}, and the routes of the literature use no other theorem of
this paper.
\end{remark}

\subsubsection*{Two of the three are not reductions}

A reader might hope that (F2) and (F3) are reductions in the ordinary sense,
statements to the effect that every Hodge class is carried, by some
correspondence, to a class on a Weil type variety or on an abelian variety,
so that the whole conjecture would follow from the Weil classes alone. Both
hopes are refutable, and we refute them, because a frontier that contains a
false statement is not a frontier.

\begin{proposition}[The Hodge ring is not generated by divisor and Weil
classes]\label{prop:mumfordproduct}
Let $X$ be a complex abelian fourfold whose Hodge group is a $\QQ$-form of an
almost direct product of three copies of $\mathrm{SL}_{2}$, acting on
$V=H^{1}(X,\QQ)$ through the tensor product $V_{1}\otimes V_{2}\otimes V_{3}$
of the three standard representations, with $\End^{0}(X)=\QQ$; such
fourfolds exist and form families of dimension one, by Mumford
\cite{Mum69}, and by \cite[2.5(i)]{MZ99} they are exactly the fourfolds
with $\End^{0}(X)=\QQ$ and Hodge group smaller than $\Sp_{8}$. Then:
\begin{enumerate}[label=\textup{(\roman*)},nosep,leftmargin=2.6em]
\item the Hodge ring of $X$ is generated by the polarisation:
  $\dim_{\QQ}\Hdg^{1}(X)=\dim_{\QQ}\Hdg^{2}(X)=1$;
\item $\dim_{\QQ}\Hdg^{1}(X\times X)=3$ and $\dim_{\QQ}\Hdg^{2}(X\times X)=8$,
  while the products of divisor classes span a subspace of dimension $6$
  in $H^{4}(X\times X,\QQ)$;
\item the two remaining classes lie in $H^{4}(X\times X,\QQ)$ and are Weil
  classes of no CM field acting on $X\times X$, on a quotient of it, or on an
  abelian subvariety of it.
\end{enumerate}
In particular the statement that every Hodge class on an abelian variety
lies in the subring generated by divisor classes and Weil classes of
quotients is false in dimension eight, though true through dimension five
\cite{MZ99}.
\end{proposition}

\begin{proof}
Hodge classes are the invariants of the Hodge group, so
$\Hdg^{k}(Y)\otimes\CC$ is the space of $\Hg(Y)_{\CC}$-invariants in
$\bigwedge^{2k}H^{1}(Y,\CC)$, and $\Hg(X\times X)=\Hg(X)$ acting
diagonally on $V\oplus V$. Over $\CC$ the group is
$\mathrm{SL}_{2}^{3}$ and the weights of $V$ under its maximal torus are the
eight sign vectors $(\pm1,\pm1,\pm1)$. Item (XXXV) of \Cref{sec:verification}
decomposes every $\bigwedge^{k}V$ by peeling highest weights and finds
\begin{gather*}
  \textstyle\bigwedge^{2}V=V(2,2,0)\oplus V(2,0,2)\oplus V(0,2,2)\oplus\mathbf{1},\\
  \textstyle\bigwedge^{4}V=V(2,2,2)\oplus\bigoplus_{\text{cyc}}V(4,0,0)
  \oplus\bigoplus_{\text{cyc}}V(2,2,0)\oplus\mathbf{1},
\end{gather*}
each with one trivial summand, which is (i). For (ii),
$\bigwedge^{4}(V\oplus V)=\bigoplus_{a}\bigwedge^{a}V\otimes\bigwedge^{4-a}V$
and the invariants in $\bigwedge^{a}V\otimes\bigwedge^{4-a}V$ are
$\Hom_{\mathrm{SL}_{2}^{3}}(\bigwedge^{a}V,\bigwedge^{4-a}V)$, since every
summand is self-dual; the counts are $1,1,4,1,1$, total $8$, the $4$ coming
from the four distinct irreducible summands of $\bigwedge^{2}V$. Likewise the
invariants in $\bigwedge^{2}(V\oplus V)$ number $1+1+1=3$. The divisor
classes of $X\times X$ are the three symplectic forms $\psi_{11}$, $\psi_{12}$,
$\psi_{22}$ on $V\oplus V$ built from the polarisation $\psi$ of $V$, and the
same item multiplies them in the exterior algebra of $V\oplus V$ and finds
their six pairwise products linearly independent; it also counts the
$\Sp_{8}$-invariants of $\bigwedge^{4}(V\oplus V)$ by the Weyl character
formula and finds six, so the six products are all the classes a Hodge group
equal to $\Sp_{8}$ would allow. For (iii), the Weil classes of an abelian
variety $Y$ with respect to a CM field $F\subseteq\End^{0}(Y)$ are, by
definition, the classes in $\bigwedge^{2k}_{F}H^{1}(Y,\QQ)\subseteq
H^{2k}(Y,\QQ)$ with $2k=\dim_{F}H^{1}(Y,\QQ)$ (\Cref{setup:cmfield}). A CM
field acting on $X\times X$ lies in $\End^{0}(X\times X)=M_{2}(\QQ)$, so it is
imaginary quadratic, $2k=8$, and its Weil classes lie in $H^{8}$, while the two
classes of (ii) have degree four. Every proper
abelian subvariety and every proper quotient of $X\times X$ is isogenous to
$X$, because $X$ is simple, and $\End^{0}(X)=\QQ$ contains no CM field, so $X$
carries no Weil classes at all. This is the statement of \cite[0.1(iii)]{MZ99}
that the exceptional classes of $X\times X$ are not of Weil type. The last
sentence is (ii) and (iii) together with the theorem of \cite{MZ99} for
dimension at most five.
\end{proof}

\begin{proposition}[No Hodge-theoretic reduction to abelian varieties]
\label{prop:noabeliantype}
Let $Y$ be a very general smooth quintic threefold in $\PP^{4}$, and let
$V=H^{3}(Y,\QQ)$ with its polarised Hodge structure of weight three. Then $V$
is not isomorphic to a subquotient of any Hodge structure of the form
$\bigoplus_{i}T_{i}$, where each $T_{i}$ is obtained from the $H^{1}$ of
finitely many abelian varieties by tensor products, duals and Tate twists. In
particular the Hodge structure $H^{*}(Y,\QQ)$ is not of abelian type, and no
correspondence between $Y$ and an abelian variety, or a product of abelian
varieties, induces an injection of $V$ into the cohomology of the latter.
\end{proposition}

\begin{proof}
Write $\mathrm{MT}(W)$ for the Mumford--Tate group of a rational Hodge
structure $W$ and $\mu_{W}\colon\mathbb{G}_{m}\to\mathrm{MT}(W)_{\CC}$ for
its Hodge cocharacter, acting on $W^{p,q}$ by $z^{p}$. The adjoint action of
$\mu_{W}$ on $\operatorname{Lie}\mathrm{MT}(W)_{\CC}\subseteq\End(W_{\CC})$ has weights
among the differences $p-p'$ of two Hodge indices of $W$.

If $W$ has weight one, these differences lie in $\{-1,0,1\}$. If $W'$ is a
subquotient of a tensor construction $T$ on $W$, then $\mathrm{MT}(W')$ is the
image of $\mathrm{MT}(W)$ acting on $W'$, and $\mu_{W'}$ is the image of
$\mu_{W}$; so $\operatorname{Lie}\mathrm{MT}(W')$ is a quotient of $\operatorname{Lie}\mathrm{MT}(W)$
as a representation of $\mathbb{G}_{m}$ through $\mu$, and its weights are a
subset of $\{-1,0,1\}$. The same holds with $W=\bigoplus H^{1}(A_{i})$.
Hence for every Hodge structure of abelian type the adjoint weights of the
Hodge cocharacter lie in $\{-1,0,1\}$.

For $V=H^{3}(Y,\QQ)$ we show that the adjoint weight $3$ occurs. The
canonical bundle of a quintic threefold is trivial by adjunction, so
$h^{3,0}(Y)=1$; pick $0\ne e\in H^{3,0}$. The intersection form $\psi$ on $V$
is symplectic and pairs $H^{3,0}$ with $H^{0,3}$ perfectly, so the rank one
endomorphism $x\mapsto\psi(x,e)\,e$ of $V_{\CC}$ is nonzero, sends $H^{0,3}$
to $H^{3,0}$ and kills the rest, and lies in $\mathfrak{sp}(V_{\CC},\psi)$
because $\psi(\psi(x,e)e,y)+\psi(x,\psi(y,e)e)
=\psi(x,e)\bigl(\psi(e,y)+\psi(y,e)\bigr)=0$; its adjoint weight is
$3-0=3$. It remains to see that $\mathfrak{sp}(V_{\CC},\psi)\subseteq
\operatorname{Lie}\mathrm{MT}(V)_{\CC}$. Let $S$ be the parameter space of smooth quintic
threefolds, connected and nonsingular, and $\mathbb{V}$ the polarisable
variation of Hodge structure $R^{3}$ of the universal family. The connected
algebraic monodromy group $G^{0}$ of $\mathbb{V}$ contains the connected
monodromy group of any Lefschetz pencil of quintics, that is of hyperplane
sections of the fifth Veronese image of $\PP^{4}$; for a Lefschetz pencil of
odd fibre dimension the image of the monodromy representation on the
$\ell$-adic vanishing cohomology $E\otimes\QQ_{\ell}$ is open in
$\Sp(E\otimes\QQ_{\ell},\psi)$ \cite[Theorem 5.10]{Del74}, so the Zariski
closure of the monodromy group over $\QQ$ is $\Sp(E,\psi)$; here
$E=H^{3}(Y_{t},\QQ)$, the vanishing cohomology being the orthogonal
complement of the image of $H^{3}(\PP^{4},\QQ)=0$. Hence $G^{0}\supseteq
\Sp(V,\psi)$, and $G^{0}\subseteq\Sp(V,\psi)$ since the monodromy preserves
$\psi$. At a very general
point $Y$ the group $G^{0}$ is a normal subgroup of the derived group of
$\mathrm{MT}(V)$ \cite[Theorem 1]{And92}, so $\Sp(V,\psi)\subseteq
\mathrm{MT}(V)^{\mathrm{der}}\subseteq\Sp(V,\psi)$, and the adjoint weight
$3$ occurs in $\operatorname{Lie}\mathrm{MT}(V)$. So $V$ is not of abelian type, and
neither is any Hodge structure containing it as a subquotient, in particular
$H^{*}(Y,\QQ)$. A correspondence inducing an injection of $V$ into the
cohomology of an abelian variety would exhibit $V$ as a sub-Hodge structure
of a Tate twist of that cohomology, which is of abelian type. Item (XXXV)
records the Hodge numbers $(1,101,101,1)$ from the Jacobian ring of the
Fermat quintic, checks the symplectic identity and the weight of the
displayed endomorphism in a four-dimensional model with the same weights, and
lists the adjoint weights that occur.
\end{proof}

\begin{remark}[What the two propositions do to the frontier]
\label{rem:notreductions}
Neither proposition changes \Cref{thm:closure}: the rule set uses (F2) and
(F3) as implications, and an implication whose hypothesis is the conjecture
for one class of Hodge classes and whose conclusion is the conjecture for a
larger class is true whatever the classes are. What does change it is
\Cref{prop:f3ishc}: the varieties of (F3) are not a class disjoint from the
abelian ones but contain them, through products with $\PP^{1}$, so (F3) is
the whole conjecture and not its non-abelian half. What they change is how the
two statements may be read. Before them one could hope that (F2) and (F3) were
structural facts awaiting proof, so that the conjecture would be a theorem
about Weil classes and nothing else; \Cref{prop:mumfordproduct} shows that the
Hodge classes on abelian varieties are not exhausted by divisor and Weil
classes once the dimension reaches eight, and \Cref{prop:noabeliantype} that
the cohomology of a general variety cannot be moved into that of abelian
varieties at all. Each of the two is therefore the Hodge conjecture for a class
of varieties on which the methods of this paper, which are methods about Weil
lines, have no purchase. The product $X\times X$ is itself of
Weil type for every imaginary quadratic $K$, through an embedding
$K\subset M_{2}(\QQ)=\End^{0}(X\times X)$ whose Rosati involution for the
product polarisation is conjugation, with signature $(4,4)$; its Weil classes
lie in $H^{8}$, and its Hodge group is far smaller than $\SU(4,4)$. The two
classes in $H^{4}$ are of a different kind, and they show that even on a
variety of Weil type the Weil line need not be the whole of the exceptional
cohomology once the Hodge group is small. They are the smallest instance of
(F2), and \Cref{ssec:mumford} identifies them and says what one K3 surface
would have to supply for them to be algebraic.
\end{remark}

\begin{figure}[tbp]
\centering
  \includegraphics[width=\textwidth]{fig_closure}
\caption{The route of this paper in \Cref{thm:closure}, drawn as a stack of
plates, one per level of the derivation; the other routes of the literature
are drawn in \Cref{fig:frontier}. The spine in indigo runs from the Weil families at the bottom to
the conjecture at the top. The two statements (F2) and (F3) enter in clay from
the right, each at the level where it is needed; (F2) is the conjecture for a
class of Hodge classes that no Weil line reaches, not a reduction, and (F3) is
the whole conjecture (\Cref{prop:f3ishc}), so the route drawn is minimal
only with (F3) replaced by \textup{(F3$'$)} of \Cref{def:f3prime}. From
below,
the propagation statement (P2) in grass reaches the two discriminant cases of
the quadratic fields and the fields of higher degree alike, because
\Cref{thm:everyfamily} and \Cref{thm:cmbasepoint} put a base point in every
family; it is read in the corrected form of \Cref{rem:p2prime}. The two
edges inside the levels are \Cref{prop:descent}: descent from the trivial
discriminant to the others, and scalar extension from the fields of degree at
least four to the quadratic fields; with them (P2) is needed only for the
split families of the fields of higher degree. The three
pillars under the grass node are the results of \Cref{thm:p2numerical},
\Cref{thm:p2support} and \Cref{cor:hhfactor} on the first form, in which the
Chern character is exactly a Weil class; by \Cref{thm:p2false} that form is
never met, so the number on the first pillar is attained by no complex. The
secant route in clay reaches the trivial discriminant, since the member it is
built on is split, and through descent the rest of the quadratic case; it has
no edge to the fields of higher degree. The bounded criterion (P1) is drawn apart, in slate, closed on
itself: it is equivalent to its own conclusion.}
\label{fig:closure}
\end{figure}

\begin{remark}[Where this leaves the secant route]\label{rem:secantplace}
\Cref{thm:closure}(iv) is a statement about the shape of the argument, not
about its difficulty, and we say plainly what it does and does not
mean.

It does not mean the secant route is worthless. It is the only route on which
anything above dimension four has ever been proved, it is the route that
settles $n=2$ at every discriminant and $n=3$ at trivial discriminant, and the
reduction of
\Cref{thm:factor} turns its remaining demand into one identity in the Yoneda
algebra of a sheaf on a fourfold, which is the form in which a demand can be
attacked. What \Cref{cor:cmdichotomy}, \Cref{thm:noci} and
\Cref{thm:nosplit} do to that identity is narrow it to two shapes of
support, neither resolved by sums of line bundles, and, in the smooth shape,
to four discriminants.

What it does mean is that finishing the secant route would not finish the
problem, and would not reduce it either. Through \Cref{prop:descent} it would
give every imaginary quadratic family, the trivial discriminant directly and
the others by descent, which is a great deal; an earlier version of this
remark said that it stops at the trivial discriminant, and that was an
omission of the rule set. It would give nothing over the CM fields of degree
at least four, whose split families need (P2) for themselves, and by
\Cref{prop:descent}(ii) that statement gives the imaginary quadratic fields as
well. Markman's extension of the secant construction to those split families
\cite{Mar25c} is aimed at exactly that statement; its semiregularity step is
open, and it is not among the rules of \Cref{setup:rules}. So the statement that would have to be proved anyway is the
propagation statement for those split families, and once proved it subsumes
the secant route entirely.

It is tempting to read (C1) and (C2) as what remains for the imaginary
quadratic case, and the computation shows why that reading is wrong: it puts
the demand for a construction where a theorem about propagation belongs.
(C1) and (C2) are demands whose satisfaction would be a considerable
achievement, would settle the Weil classes of every imaginary quadratic field,
and would close nothing that (P2) does not close.
\end{remark}


\subsection{Summary}

\begingroup
\renewcommand{\arraystretch}{1.25}
\begin{longtable}{@{}>{\raggedright\arraybackslash}p{0.30\linewidth}>{\raggedright\arraybackslash}p{0.30\linewidth}>{\raggedright\arraybackslash}p{0.30\linewidth}@{}}
\caption{The routes considered, and the boundary of each result.}
\label{tab:scope}\\
\toprule
\textbf{Route} & \textbf{Status here} & \textbf{Not covered} \\
\midrule
\endfirsthead
\multicolumn{3}{@{}l}{\footnotesize\itshape Table \thetable, continued.}\\[2pt]
\toprule
\textbf{Route} & \textbf{Status here} & \textbf{Not covered} \\
\midrule
\endhead
\midrule
\multicolumn{3}{r@{}}{\footnotesize\itshape continued on the next page}\\
\endfoot
\bottomrule
\endlastfoot
Ambient object, restricted to $\Av$, transported by $\Phi_{\cP}$
  & closed by \Cref{thm:divisor}
  & ambient varieties with exceptional classes \\
Any $K$-equivariant construction with the norm character
  & closed by \Cref{thm:character}
  & constructions that are not eigenclasses \\
Secant dimension count
  & no obstruction, \Cref{cor:nonumerical}
  & none \\
Semiregularity dimension count
  & no obstruction, \Cref{cor:countnotbind}
  & none \\
Semiregularity of an object whose Chern character sits on the Weil line
  & closed by \Cref{thm:rankzero} and \Cref{cor:kernelfamily} for direct
    sums, and for every object by \Cref{thm:p2false}
  & objects of positive rank; the weaker hypothesis of
    \cite[Question 11.4]{Mar25b} \\
Induction through hyperplane sections
  & closed by \Cref{thm:hyperplane}, \Cref{thm:returntrip}
  & correspondences, degenerations, coniveau \\
Kuga--Satake, from a weight two structure of K3 type
  & not obstructed, \Cref{prop:ksnotrestriction}
  & available only at $n=1$, \Cref{cor:ksboundary} \\
Divisors, on a member with quaternionic multiplication
  & succeeds, \Cref{thm:quatweil}
  & fails at full Hodge group, \Cref{prop:divfails} \\
Propagation off the base locus
  & reduced to one bound or one rank, \Cref{thm:degreebound},
    \Cref{thm:semiregclosure}
  & neither hypothesis proved, \Cref{rem:tworemain} \\
Objects with effective Chern character at a base point
  & closed by \Cref{thm:positivity}
  & objects indecomposable in $D^{b}$, \Cref{rem:whatitmustbe} \\
The weakened criterion for a secant object of positive rank
  & reduced to one identity on the factor, \Cref{thm:factor}; the identity
    fails for Markman's candidate, \Cref{thm:candidatefails}, and at every
    transversal double point, \Cref{cor:localobstruction}
  & objects whose support has no such point, \Cref{rem:candidatemeaning} \\
A secant object of rank one with smooth support
  & invariants forced and the discriminant bounded by $9$,
    \Cref{thm:smoothinvariants}, \Cref{cor:fourdiscriminants}
  & existence for $d\in\{1,3,5,7\}$, and one vanishing in
    $H^{1}(S,N_{S/X})$, \Cref{rem:smoothremains} \\
The local part of the weakened criterion
  & it is a condition on depth: no obstruction at a Cohen--Macaulay point,
    \Cref{thm:cmvanishing}, \Cref{cor:cmdichotomy}
  & Cohen--Macaulay supports, where the criterion becomes one global
    vanishing \\
A support that is a complete intersection
  & closed for every $b$ and every $d$, \Cref{thm:noci}
  & supports of higher Hilbert--Burch rank, \Cref{thm:burchsearch} \\
A support with a Hilbert--Burch resolution by line bundles or by simple
semi-homogeneous bundles
  & closed for every rank, twist and $d$, \Cref{thm:nosplit},
    \Cref{thm:noblocks}
  & resolutions over bundles that are not projectively flat \\
The secant route, granted every demand it makes
  & reaches $\delta=\delta_{0}$, and every $\delta$ by \Cref{prop:descent},
    \Cref{thm:closure}(iv)
  & the CM fields of degree at least four, which give the quadratic ones by
    \Cref{prop:descent}(ii) \\
Propagation from a base point
  & closes every imaginary quadratic family at once, \Cref{thm:degreebound},
    \Cref{thm:semiregclosure}
  & (P2) not proved, \Cref{cor:fourremain} \\
The bounded form (P1), read on the transported cycle
  & closed: the transport spreads each component by $c^{2n}/m_{i}$ and the
    data is unbounded whenever a component has finite stabiliser,
    \Cref{lem:multiplicity}, \Cref{thm:p1forces}
  & nothing; the other reading is closed too \\
The bounded form (P1), read on the minimal representative
  & closed: equivalent to the conjecture for the family,
    \Cref{thm:p1equivalent}
  & nothing; it is a restatement, so it is not a route \\
The semiregularity form (P2), for a given representative
  & forces the evaluation map to vanish on the $4n^{2}$-dimensional
    annihilator, and excludes every representative that is a sheaf on a
    smooth germ of codimension two of its support, \Cref{thm:p2forces},
    \Cref{cor:p2shape}
  & representatives that are genuine complexes along every such component;
    members with an abelian subvariety tangent to an eigenspace \\
The same form, tested by a number
  & $\dim\Ext^{2}(E,E)=2n(2n-1)$ at one base point would suffice,
    \Cref{thm:p2numerical}, and no representative attains it,
    \Cref{thm:p2false}
  & nothing \\
The same form, for a representative in codimension $n$
  & closed at every member of every family: the support must have a
    component of codimension less than $n$, and at a general member every
    component of codimension at least two has multiplicity zero,
    \Cref{lem:finiteproduct}, \Cref{thm:p2support}
  & representatives glued along components of lower codimension with
    alternating length zero \\
The Hochschild action on such a representative
  & the annihilator is the mixed part of a splitting of $HH^{1}$ and the
    action factors through a ring of dimension $2^{2n+1}-1$ in one of
    dimension $2^{4n}$, \Cref{thm:hhsplitting}, \Cref{cor:hhfactor}
  & whether a complex with that action exists \\
The shape of such a representative
  & an indecomposable summand carries the class; at $n=2$, and at $n=3$
    granting the compatibility of \Cref{cor:hhfactor}(iii), its
    endomorphism algebra has dimension at least three, so no simple object
    serves, \Cref{thm:p2endo}
  & nothing: no object of any shape meets the pure form, \Cref{thm:p2false} \\
The pure form of (P2), Chern character exactly a Weil class
  & false for every imaginary quadratic field and every $n\ge2$: a
    semiregular such complex would deform to very general non-algebraic Weil
    tori, whose coherent sheaves have no Chern character below the top
    degree, \Cref{prop:weiltorussheaves}, \Cref{thm:p2false}
  & the corrected form (P2), with powers of the polarisation allowed in the
    Chern character, \Cref{rem:p2prime} \\
The reduction to Weil classes, and to abelian varieties
  & false as a statement about Hodge rings in dimension eight,
    \Cref{prop:mumfordproduct}; impossible Hodge-theoretically,
    \Cref{prop:noabeliantype}
  & the algebraicity of the classes beyond the Weil lines, which is (F2), and
    (F3), which is the whole conjecture, \Cref{prop:f3ishc} \\
The two classes on the self-product of a Mumford fourfold
  & a real multiplication by a cubic field on the primitive part of $H^{2}$,
    algebraic once the Kuga--Satake class of one K3 surface of Picard number
    thirteen is, \Cref{prop:mumfordrm}, \Cref{thm:mumfordks}
  & that Kuga--Satake class, \Cref{rem:mumfordks} \\
Bypasses for the same two classes
  & degeneration and larger families closed off by rigidity; semiregularity
    reduced to $\dim\Ext^{2}=119$ at one CM point; specialisation reduced to
    a degree bound, \Cref{thm:mumfordrigid}, \Cref{thm:mumfordnumerical},
    \Cref{cor:mumfordbounded}
  & an object with $\dim\Ext^{2}=119$, or the degree bound,
    \Cref{rem:bypassaudit} \\
Objects built from line bundles, for the same two classes
  & closed at every point of the curve: the Chern characters stay invariant
    under the Lefschetz group, \Cref{thm:lefschetzclosure}
  & steps that break that symmetry, \Cref{rem:leavel} \\
Twistor deformation of the square, for the same two classes
  & closed: the component of the Hodge locus among all complex tori is the
    Mumford curve, which contains no twistor line, \Cref{prop:notwistor}
  & the Kuga--Satake class, which moves in a seven-dimensional family \\
The five routes left open, for the same two classes
  & one statement: three of them are equivalent to the target and two imply
    it, and every minimal set of open statements yielding the target has one
    element, \Cref{cor:mumfordone}, item (XLIX)
  & the target itself, which is only the first case of (F2),
    \Cref{rem:mumfordneeded} \\
Pull-back from a hyperk\"ahler manifold, for the same two classes
  & the square is holomorphic symplectic and the classes are the twisted dual
    forms of the part of $H^{2}$ its symplectic class generates; a symplectic
    rational map to some $S_{\mu}^{[n]}$, $n\ge5$, suffices, and no K3
    surface or hyperk\"ahler eightfold serves, \Cref{prop:hksquare}
  & such a map, \Cref{rem:hkroute} \\
Reduction modulo $p$, for the same two classes
  & represented by cycles at every closed point of every reduction; the
    target is equivalent to a bound on their Hilbert data on a Zariski dense
    set, \Cref{prop:modp}, \Cref{thm:modpdensity}
  & the bound, \Cref{rem:modpremains} \\
The Weil classes of every CM field
  & the same propagation problem, with a CM base point in every family,
    \Cref{thm:cmbasepoint}, \Cref{thm:cmpropagation}
  & (P2) not proved for any field \\
What is left, as a whole
  & five minimal sets of open statements: $\{\mathrm{F3}\}$ and
    $\{\mathrm{L},\mathrm{M}\}$, each equivalent to the conjecture, and
    three through the conjecture modulo abelian varieties,
    $\{\mathrm{L},\mathrm{F3}'\}$, $\{\mathrm{V},\mathrm{F3}'\}$ and this
    route, $\{\mathrm{P2}_{\mathrm{split}},\mathrm{F2},\mathrm{F3}'\}$,
    \Cref{thm:closure}, \Cref{prop:f3ishc}, \Cref{cor:fourremain},
    \Cref{rem:bullseye}
  & the Lefschetz standard conjecture, (M), (V), (F2), (F3$'$), all open \\
The conjecture modulo abelian varieties, (F3$'$)
  & holds on the class $\mathcal{A}$ of varieties whose cohomology is
    reached from abelian varieties by algebraic correspondences (curves,
    abelian varieties, products, surjective images), in degrees $0$, $2$,
    $2n-2$, $2n$, and in dimension at most three, \Cref{prop:f3prime}
  & degree four on fourfolds outside $\mathcal{A}$ and beyond, for
    instance squares of K3 surfaces without an algebraic Kuga--Satake
    correspondence, \Cref{rem:f3primeopen} \\
\end{longtable}
\endgroup
